\chapter{Teoria}
\label{cap:teoria}

\section{Anell dels polinomis}

S'anomena \textbf{polinomi en la indeterminada} (o variable) $x$ tota
expressió reduïble a la forma:
\begin{equation*}
a_{0}+a_{1}x+a_{2}x^{2}+\cdots +a_{n-1}x^{n-1}+a_{n}x^{n}
\end{equation*}
on $n$ és un nombre natural i $a_{0},a_{1},a_{2},\ldots,a_{n-1},a_{n}$
són nombres reals qualssevol. La lletra $x$ que apareix en el
polinomi no designa res en particular i, per això, s’anomena
indeterminada; en realitat, pot substituir-se per qualsevol altra lletra o
signe. Si $a_{n}\neq 0$, aleshores es diu que $n$ és el \textbf{grau }del
polinomi; s'admet que un nombre real diferent de zero és un polinomi
de grau zero i que el nombre real $0$ és un polinomi sense grau, s'anomena \textbf{polinomi zero}. A cadascun dels sumands $a_{i}x^{i}$ ($i=1,2,...,n$) del polinomi l'anomenem \textbf{monomi de grau} $i$ i
aleshores $a_{i}$ s’anomena \textbf{coeficient del monomi de grau} $i$. El terme $a_{0}$ s'anomena \textbf{terme independent} del polinomi (o
coeficient del monomi de grau zero si $a_{0}\neq 0$).

En general, els nombres reals $a_{0},a_{1},\ldots,a_{n-1},a_{n}$ s'anomenen també \textbf{coeficients del polinomi} i, per a referir-nos a un d'ells, direm també que $a_{i}$ és el \textbf{coeficient del terme de grau} $i$ del polinomi. S'anomena \textbf{nombre de termes} del polinomi el nombre de monomis no nuls que apareixen en el
polinomi; si són dos, el polinomi s’anomena \textbf{binomi}, i si són
tres, \textbf{trinomi}. Finalment, per a referir-nos a un polinomi en la
indeterminada $x$ farem servir la notació $p(x)$. D'aquesta manera, si $p(x)$ és un polinomi de grau $n$ amb coeficients $a_{0},a_{1},\ldots,a_{n-1},a_{n}$, escriurem
\begin{equation*}
p(x)=a_{0}+a_{1}x+a_{2}x^{2}+\cdots +a_{n-1}x^{n-1}+a_{n}x^{n}
\end{equation*}
Dos polinomis són \textbf{iguals} si coincideixen els coeficients de tots
els termes del mateix grau en un i l'altre polinomi. És clar que si
dos polinomis són iguals, aleshores tots dos tenen el mateix grau; en
particular, donats dos polinomis de grau $n$
\begin{equation*}
p(x)=a_{0}+a_{1}x+a_{2}x^{2}+\cdots +a_{n-1}x^{n-1}+a_{n}x^{n}~\ \
(a_{n}\neq 0)
\end{equation*}i\begin{equation*}
q(x)=b_{0}+b_{1}x+b_{2}x^{2}+\cdots +b_{n-1}x^{n-1}+b_{n}x^{n}~\ \
(b_{n}\neq 0)
\end{equation*}direm que $p(x)=q(x)$ si
\begin{equation*}
a_{k}=b_{k}~\ \ \ \ (k=0,1,2,...,n)
\end{equation*}

Per exemple,\begin{equation*}
p(x)=-1+2x-3x^{2}+5x^{4}
\end{equation*}és un polinomi de grau $4$ amb coeficients $-1,2,-3$ i $5$, essent el seu terme independent $-1$. Observa que $p(x)$ està format per una
suma de monomis: $-1$ és un monomi de grau zero, $2x$ és un monomi de
grau $1$, $-3x^{2}$ és un monomi de grau $2$ i $5x^{4}$ és un monomi de
grau $4$. Per exemple, el coeficient del monomi de grau $2$ és $-3$ i el coeficient
del monomi de grau $3$ és $0$

De vegades un polinomi no és donat en forma reduïda. Això passa, per
exemple, amb el següent polinomi\begin{equation*}
p(x)=x^{3}-3x^{2}-\sqrt{2}+\sqrt{3}x-7x^{2}+5x^{3}+1-6x^{3}+10x^{2}+x
\end{equation*}Aleshores, efectuant operacions amb els monomis del mateix grau,
anomenats \textbf{termes semblants}, obtenim\begin{eqnarray*}
p(x) &=&-\sqrt{2}+1+(\sqrt{3}+1)x+(-3-7+10)x^{2}+(1+5-6)x^{3} \\
&=&-\sqrt{2}+1+(\sqrt{3}+1)x
\end{eqnarray*}Mitjançant aquesta \textbf{reducció de termes semblants} veiem que
resulta un polinomi de grau $1$ amb coeficients $-\sqrt{2}+1$ i $\sqrt{3}+1$.

\subsection{Valor numèric d'un polinomi}

Hem dit que la indeterminada d'un polinomi pot ser qualsevol cosa.
En particular, un polinomi\begin{equation*}
p(x)=a_{0}+a_{1}x+a_{2}x^{2}+\cdots +a_{n-1}x^{n-1}+a_{n}x^{n}
\end{equation*}pren un valor quan a la indeterminada $x$ li assignem un nombre real.
S'anomena \textbf{valor numèric} d'un polinomi $p(x)$ per a $x=k$ ($k$
és un nombre real qualsevol) el nombre\begin{equation*}
p(k)=a_{0}+a_{1}k+a_{2}k^{2}+\cdots +a_{n-1}k^{n-1}+a_{n}k^{n}
\end{equation*}és a dir, al resultat de substituir la $x$ pel nombre $k$ i efectuar
després les operacions.

Per exemple, el valor numèric del polinomi $p(x)=1+x-x^{3}$ per a $x=-1$
és\begin{eqnarray*}
p(-1) &=&1+(-1)-(-1)^{3} \\
&=&1-1+1 \\
&=&1
\end{eqnarray*}

\subsection{Suma de polinomis}

Donats dos polinomis\begin{equation*}
\begin{array}{l}
p(x)=a_{0}+a_{1}x+a_{2}x^{2}+\cdots +a_{n-1}x^{n-1}+a_{n}x^{n} \\
q(x)=b_{0}+b_{1}x+b_{2}x^{2}+\cdots +b_{m-1}x^{m-1}+b_{m}x^{m}\end{array}\end{equation*}amb $a_{n}\neq 0$ i $b_{m}\neq 0$, o sigui, de graus $n$ i $m$
respectivament. Definim la \textbf{suma de dos monomis del mateix grau} $i$ de la manera següent\begin{equation*}
a_{i}x^{i}+b_{i}x^{i}=(a_{i}+b_{i})x^{i}
\end{equation*}
Observa que, quan $x$ pren algun valor numèric, aquesta definició és coherent amb el fet que el producte de nombres reals
sigui distributiu respecte de la suma. Aleshores, definim \textbf{polinomi
suma} el polinomi $p(x)+q(x)$ que resulta de sumar els monomis del mateix
grau. Així, tenim\begin{equation*}
p(x)+q(x)=(a_{0}+b_{0})+(a_{1}+b_{1})x+(a_{2}+b_{2})x^{2}+\cdots
\end{equation*}Si en un polinomi o en tots dos falta algun monomi, s'ha de suposar que
el seu coeficient és zero.

Per exemple, si\begin{equation*}
\begin{array}{l}
p(x)=2-3x+5x^{2} \\
q(x)=-1+3x^{2}-2x^{3}\end{array}\end{equation*}aleshores el polinomi suma $p(x)+q(x)$ és\begin{equation*}
\begin{array}{cccccc}
p(x) & = & 2 & -3x & 5x^{2} & 0x^{3} \\
q(x) & = & -1 & 0x & 3x^{2} & -2x^{3} \\ \hline
p(x)+q(x) & = & 1 & -3x & 8x^{2} & -2x^{3}\end{array}\end{equation*}és a dir,\begin{equation*}
p(x)+q(x)=1-3x+8x^{2}-2x^{3}
\end{equation*}

\subsubsection{Propietats de la suma}

\begin{enumerate}
\item[(1)] \textbf{Associativa}: qualssevol que siguin els polinomis $p(x),q(x)$ i $r(x)$ es compleix\begin{equation*}
\left[ p(x)+q(x)\right] +r(x)=p(x)+\left[ q(x)+r(x)\right]
\end{equation*}

\item[(2)] \textbf{Commutativa}: qualssevol que siguin els polinomis $p(x)$
i $q(x)$ es compleix\begin{equation*}
p(x)+q(x)=q(x)+p(x)
\end{equation*}

\item[(3)] \textbf{Existència d'element neutre}: per a qualsevol polinomi $p(x)$ es compleix\begin{equation*}
p(x)+0=p(x)
\end{equation*}essent $0$ el polinomi zero, és a dir, el polinomi que té tots les seves
coeficients iguals a $0$.

\item[(4)] \textbf{Existència d'element oposat}: per a cada polinomi $p(x) $existeix $-p(x)$, anomenat \textbf{polinomi oposat} de $p(x)$, tal
que\begin{equation*}
p(x)+\left[ -p(x)\right] =0
\end{equation*}Aquest polinomi és únic per a cada $p(x)$ i, si $p(x)=a_{0}+a_{1}x+a_{2}x^{2}+\cdots +a_{n-1}x^{n-1}+a_{n}x^{n}$, aleshores\begin{equation*}
-p(x)=(-a_{0})+(-a_{1})x+(-a_{2})x^{2}+\cdots
+(-a_{n-1})x^{n-1}+(-a_{n})x^{n}
\end{equation*}és a dir, $-p(x)$ té per coeficients els oposats dels coeficients de
$p(x)$. Com és habitual, definim\begin{equation*}
p(x)-q(x)=p(x)+\left[ -q(x)\right]
\end{equation*}és a dir, la \textbf{resta} de dos polinomis $p(x)-q(x)$ és la suma de $p(x) $amb el oposat de $q(x)$.

\item[(5)] \textbf{Grau del polinomi suma}: el grau del polinomi suma és
més petit o igual que el màxim dels graus dels
polinomis sumands\begin{equation*}
\text{grad\,}\left[ p(x)+q(x)\right] \leq \max \left\{ \text{grad
\,}p(x),\text{grad\,}q(x)\right\}
\end{equation*}Per exemple, si\begin{equation*}
\begin{array}{l}
p(x)=-16+x^{2}-6x^{4}+x^{5} \\
q(x)=1+2x^{2}-5x^{3}\end{array}\end{equation*}aleshores\begin{equation*}
p(x)+q(x)=-15+3x^{2}-5x^{3}-6x^{4}+x^{5}
\end{equation*}i es compleix\begin{equation*}
\text{grad\,}\left[ p(x)+q(x)\right] =5\leq 5=\max
\left\{\text{grad\,}p(x),\text{grad\,}q(x)\right\}
\end{equation*}Però, si\begin{equation*}
\begin{array}{l}
p(x)=-16+x^{2}-6x^{4}+x^{5} \\
q(x)=1+2x^{2}-x^{5}\end{array}\end{equation*}aleshores\begin{equation*}
p(x)+q(x)=-15+3x^{2}-6x^{4}
\end{equation*}i també es compleix\begin{equation*}
\text{grad\,}\left[ p(x)+q(x)\right] =4\leq 5=\max \left\{
\text{grad\,}p(x),\text{grad\,}q(x)\right\}
\end{equation*}
\end{enumerate}

\subsection{Producte de polinomis}

Donats dos polinomis\begin{equation*}
\begin{array}{l}
p(x)=a_{0}+a_{1}x+a_{2}x^{2}+\cdots +a_{n-1}x^{n-1}+a_{n}x^{n} \\
q(x)=b_{0}+b_{1}x+b_{2}x^{2}+\cdots +b_{m-1}x^{m-1}+b_{m}x^{m}\end{array}\end{equation*}amb $a_{n}\neq 0$ i $b_{m}\neq 0$, o sigui, de graus $n$ i $m$
respectivament. Definim el \textbf{producte de dos monomis }qualssevol
de la manera següent\begin{equation*}
\left( a_{i}x^{i}\right) \cdot \left( b_{j}x^{j}\right) =\left(
a_{i}b_{j}\right) ~x^{i+j}
\end{equation*}
Observa que, quan $x$ pren algun valor numèric, aquesta definició és coherent amb les propietats commutativa i associativa del
producte de nombres reals i també amb la propietat del producte
de potències. Aleshores, definim \textbf{polinomi producte} el polinomi $p(x)\cdot q(x)$ que resulta de multiplicar cada monomi de $p(x)$ per tots
els monomis de $q(x)$ i després sumar-los tots. Així, tenim\begin{equation*}
p(x)\cdot
q(x)=(a_{0}b_{0})+(a_{0}b_{1}+a_{1}b_{0})x+(a_{0}b_{2}+a_{1}b_{1}+a_{2}b_{0})x^{2}+\cdots
\end{equation*}

Per exemple, si\begin{equation*}
\begin{array}{l}
p(x)=2-3x+5x^{2} \\
q(x)=-1+3x^{2}-2x^{3}\end{array}\end{equation*}aleshores el polinomi producte $p(x)\cdot q(x)$ és\begin{equation*}
\begin{array}{cccccccccccccc}
p(x) & = & & & & & & & & 2 & - & 3x & + & 5x^{2} \\
q(x) & = & & & & & & & - & 1 & + & 3x^{2} & - & 2x^{3} \\ \hline
& & - & 2 & & & + & 6x^{2} & - & 4x^{3} & & & & \\
& & & & & 3x & & & - & 9x^{3} & + & 6x^{4} & & \\
& & & & & & - & 5x^{2} & & & + & 15x^{4} & - & 10x^{5} \\ \hline
p(x)\cdot q(x) & = & - & 2 & + & 3x & + & x^{2} & - & 13x^{3} & + & 21x^{4}
& - & 10x^{5}\end{array}\end{equation*}

\subsubsection{Propietats del producte}

\begin{enumerate}
\item[(1)] \textbf{Associativa}: qualssevol que siguin els polinomis $p(x),q(x)$ i $r(x)$ es compleix\begin{equation*}
\left[ p(x)\cdot q(x)\right] \cdot r(x)=p(x)\cdot \left[ q(x)\cdot r(x)\right]
\end{equation*}

\item[(2)] \textbf{Commutativa}: qualssevol que siguin els polinomis $p(x)$
i $q(x)$ es compleix\begin{equation*}
p(x)\cdot q(x)=q(x)\cdot p(x)
\end{equation*}

\item[(3)] \textbf{Existència d'element unitat}: per a qualsevol polinomi $p(x)$ es compleix\begin{equation*}
p(x)\cdot 1=p(x)
\end{equation*}essent $1$ la unitat dels nombres reals, és a dir, un polinomi de
grau zero. A $1$ s'anomena aleshores \textbf{polinomi unitat}.

\item[(4)] \textbf{Distributiva del producte respecte de la suma}:
qualssevol que siguin els polinomis $p(x),q(x)$ i $r(x)$ es compleix\begin{equation*}
p(x)\cdot \left[ q(x)+r(x)\right] =p(x)\cdot q(x)+p(x)\cdot r(x)
\end{equation*}

\item[(5)] \textbf{No té divisors de zero}: si $p(x)\neq 0$ i $q(x)\neq
0 $, aleshores\begin{equation*}
p(x)\cdot q(x)\neq 0
\end{equation*}o dit d'una altra manera equivalent, si $p(x)\cdot q(x)=0$, aleshores $p(x)=0$
o $q(x)=0$.

\item[(6)] \textbf{Grau del polinomi producte}: si els polinomis
multiplicants són no nuls, aleshores el grau del polinomi producte es
igual a la suma dels graus dels polinomis\begin{equation*}
\text{grad\,}\left[ p(x)\cdot q(x)\right]
=\text{grad\,}p(x)+\text{grad\,}q(x)
\end{equation*}amb $p(x)\neq 0$ i $q(x)\neq 0$. Per exemple, si
\begin{equation*}
\begin{array}{l}
p(x)=2-3x+5x^{2} \\
q(x)=-1+3x^{2}-2x^{3}\end{array}\end{equation*}aleshores\begin{equation*}
p(x)\cdot q(x)=-2+3x+x^{2}-13x^{3}+21x^{4}-10x^{5}
\end{equation*}i es compleix\begin{equation*}
\text{grad\,}\left[ p(x)\cdot q(x)\right]
=5=2+3=\text{grad\,}p(x)+\text{grad\,}q(x)
\end{equation*}
\end{enumerate}

\subsection{Anell dels polinomis}

Si designem per $\mathbb{R}[x]$ el conjunt de tots els polinomis amb
coeficients reals, aleshores diem que $\mathbb{R}[x]$ té estructura
de \textbf{anell unitari commutatiu sense divisors de zero}. Això és així, perquè les operacions de suma i producte de polinomis són internes
(la suma de dos polinomis és un altre polinomi, i el producte de dos
polinomis és un altre polinomi) i, (i) la suma de polinomis compleix les
propietats (1) a (4) i (ii) el producte de polinomis compleix les
propietats (1) a (5). Designem per $\mathbb{Z}[x]$ el conjunt de tots
els polinomis amb coeficients enters i per $\mathbb{Q}[x]$, el conjunt
de tots els polinomis amb coeficients racionals.

Observa que $\mathbb{R}\subset \mathbb{R}[x]$, perquè, hem assumit el
conveni de considerar els nombres reals diferents de zero com
polinomis de grau zero; el nombre real $0$ és un polinomi sense grau.
Aquesta inclusió dels nombres reals en el anell dels polinomis
és tal que tota propietat relativa a aquest anell serà coherent amb
les propietats ja conegudes dels nombres reals. Per exemple, les
propietats de la suma de polinomis es tradueixen en les propietats de la
suma de nombres reals al considerar els polinomis de grau zero i el nombre real $0$.

\section{Potència d'un binomi. Fórmula de Newton}

Considerem ara el cas particular que el polinomi $p(x)$ sigui un
binomi. S'anomena potència $n$-èsima del binomi al producte de $p(x)$
per si mateix $n$ vegades, és a dir,\begin{equation*}
\left[ p(x)\right] ^{n}=p(x)\cdot \overset{n\text{ vegades}}{\cdots }\cdot
~p(x)
\end{equation*}Per a trobar una expressió general d'aquesta potència anem a suposar que
el binomi és de la forma $A+B$, essent $A$ i $B$ monomis de graus
qualssevol. A partir de les propietats de la suma i producte de monomis
així com de les propietats dels nombres combinatoris \footnote{Recorda que el nombre combinatori
\begin{equation*}
\binom{n}{m}=\frac{n!}{m!\cdot (n-m)!}
\end{equation*}és el nombre de subconjunts de $m$ elements que poden obtenir-se a
partir d'un conjunt de $n$ elements o, de forma equivalent, és el nombre de col·leccions de $m$ objectes que es poden fer a partir d'una
col·lecció de $n$ objectes, prescindint de l'ordre.}, es compleix la
\textbf{fórmula de Newton}\begin{equation*}
(A+B)^{n}=\binom{n}{0}A^{n}+\binom{n}{1}A^{n-1}B+\binom{n}{2}A^{n-2}B+\cdots
+\binom{n}{n-1}AB^{n-1}+\binom{n}{n}B^{n}
\end{equation*}En efecte, demostrarem aquesta fórmula per inducció \footnote{El mètode de demostració anomenat per inducció s'utilitza a
sovint per a provar que una determinada propietat o fórmula és vàlida
per a tot nombre natural. La prova consiste en demostrar dos fets:
(1) que la propietat o fórmula es compleix per a $n=1$ i (2) que per a tot nombre natural $k$, si la propietat es compleix per a $n=k$ (hipòtesi
d'inducció), aleshores també es compleix per al seu següent $n=k+1$.
Una vegada demostrats aquests dos fets, pel principi d'inducció
dels nombres naturals, la propietat es compleix per a tot nombre
natural $n$.}. Per a $n=1$, cal demostrar que es compleix la següent
igualtat\begin{equation*}
(A+B)^{1}=\binom{1}{0}A^{1}+\binom{1}{1}A^{0}B
\end{equation*}Però, això és evident ja que\begin{equation*}
\binom{1}{0}=\binom{1}{1}=1
\end{equation*}i per definició $A^{0}=1$ (Aquest fet és coherent amb la propietat
que $a^{0}=1$ per a tot nombre real diferent de zero). Suposem
ara que la fórmula es compleix per a $n=k$ (hipòtesi d'inducció), essent $k$ qualsevol nombre natural\begin{equation*}
(A+B)^{k}=\binom{k}{0}A^{k}+\binom{k}{1}A^{k-1}B+\binom{k}{2}A^{k-2}B+\cdots
+\binom{k}{k-1}AB^{k-1}+\binom{k}{k}B^{k}
\end{equation*}Es tracta ara de demostrar que també es compleix per a $n=k+1$. En
efecte,\begin{eqnarray*}
(A+B)^{k+1} &=&(A+B)^{k}(A+B) \\
&=&\left[ \binom{k}{0}A^{k}+\binom{k}{1}A^{k-1}B+\binom{k}{2}A^{k-2}B+\cdots
+\binom{k}{k-1}AB^{k-1}+\binom{k}{k}B^{k}\right] (A+B) \\
&=&\binom{k}{0}A^{k+1}+\binom{k}{1}A^{k}B+\binom{k}{2}A^{k-1}B+\cdots +\binom{k}{k-1}A^{2}B^{k-1}+\binom{k}{k}AB^{k}+ \\
&&\binom{k}{0}A^{k}B+\binom{k}{1}A^{k-1}B^{2}+\binom{k}{2}A^{k-2}B^{2}+\cdots +\binom{k}{k-1}AB^{k}+\binom{k}{k}B^{k+1} \\
&=&\binom{k}{0}A^{k+1}+\left[ \binom{k}{0}+\binom{k}{1}\right] A^{k}B+\left[
\binom{k}{1}+\binom{k}{2}\right] A^{k-1}B^{2}+\cdots +\left[ \binom{k}{k-2}+\binom{k}{k-1}\right] A^{2}B^{k-1}+ \\
&&\left[ \binom{k}{k-1}+\binom{k}{k}\right] AB^{k}+\binom{k}{k}B^{k+1} \\
&=&\binom{k+1}{0}A^{k+1}+\binom{k+1}{1}A^{k}B+\binom{k+1}{2}A^{k-1}B^{2}+\cdots +\binom{k+1}{k-1}A^{2}B^{k-1}+\binom{k+1}{k}AB^{k}+\binom{k+1}{k+1}B^{k+1}
\end{eqnarray*}perquè, sabem que es compleixen les següents propietats dels nombres
combinatoris\begin{equation*}
\binom{n}{0}=\binom{n}{n}=1
\end{equation*}i\begin{equation*}
\binom{n}{m}+\binom{n}{m+1}=\binom{n+1}{m+1}
\end{equation*}

Per exemple, si
\begin{equation*}
\begin{array}{c}
A=x^{2} \\
B=1\end{array}\end{equation*}aleshores\begin{eqnarray*}
(x^{2}+1)^{4} &=&\binom{4}{0}\left( x^{2}\right) ^{4}+\binom{4}{1}\left(
x^{2}\right) ^{3}\cdot 1+\binom{4}{2}\left( x^{2}\right) ^{2}\cdot 1^{2}+\binom{4}{3}\left( x^{2}\right) \cdot 1^{3}+\binom{4}{4}1^{4} \\
&=&x^{8}+4x^{6}+6x^{4}+4x^{2}+1
\end{eqnarray*}

És important observar que la fórmula de Newton no només serà vàlida per a desenvolupar expressions de la forma $(A+B)^{n}$, amb $A$ i $B$
monomis de graus qualssevol, sinó que també ho serà per a
desenvolupar expressions del mateix tipus on $A$ i $B$ siguin nombres
reals o expressions subjectes a les mateixes propietats.

\subsection{Càlcul dels coeficients en la fórmula de Newton}

Per a calcular els nombres combinatoris que apareixen en la fórmula
de Newton és útil formar el \textbf{triangle de Tartaglia}\begin{equation*}
\begin{array}{ccccccccccccc}
& & & & & 1 & & 1 & & & & & \\
& & & & 1 & & 2 & & 1 & & & & \\
& & & 1 & & 3 & & 3 & & 1 & & & \\
& & 1 & & 4 & & 6 & & 4 & & 1 & & \\
& 1 & & 5 & & 10 & & 10 & & 5 & & 1 & \\
1 & & 6 & & 15 & & 20 & & 15 & & 6 & & 1 \\
& \cdots & & \cdots & & \cdots & & \cdots & & \cdots & & \cdots &
\end{array}\end{equation*}Aquest triangle té com base la propietat dels nombres
combinatoris segons la qual\begin{equation*}
\binom{n}{m}+\binom{n}{m+1}=\binom{n+1}{m+1}
\end{equation*}a més que\begin{equation*}
\binom{n}{0}=\binom{n}{n}=1
\end{equation*}

Per exemple, per a calcular $(a+b)^{5}$ primer construirem el triangle de Tartaglia fins a la fila 5. Els nombres d'aquesta fila es
corresponen en ordre amb els nombres combinatoris que apareixen en el
desenvolupament per la fórmula de Newton de $(a+b)^{5}$. Així tenim\begin{eqnarray*}
(a+b)^{5} &=&\binom{5}{0}a^{5}+\binom{5}{1}a^{4}b+\binom{5}{2}a^{3}b^{2}+\binom{5}{3}a^{2}b^{3}+\binom{5}{4}ab^{4}+\binom{5}{5}b^{5} \\
&=&1\cdot a^{5}+5\cdot a^{4}b+10\cdot a^{3}b^{2}+10\cdot a^{2}b^{3}+5\cdot
ab^{4}+1\cdot b^{5} \\
&=&a^{5}+5a^{4}b+10a^{3}b^{2}+10a^{2}b^{3}+5ab^{4}+b^{5}
\end{eqnarray*}

\subsection{Càlcul d'un terme del desenvolupament d'un binomi}

Observa que en la fórmula de Newton\begin{equation*}
(A+B)^{n}=\binom{n}{0}A^{n}+\binom{n}{1}A^{n-1}B+\binom{n}{2}A^{n-2}B+\cdots
+\binom{n}{n-1}AB^{n-1}+\binom{n}{n}B^{n}
\end{equation*}apareixen $n+1$ termes i en cada terme hi ha un nombre
combinatori que té com numerador sempre $n$ i com denominador un nombre que és una unitat més petit que el lloc que ocupa aquest terme
com a sumand del desenvolupament. A més, prenent per definició que $A^{0}=B^{0}=1$, en cada terme apareix el producte d’una potència de $A $per una altra de $B$ de manera que l'exponent de $A$ és la diferència de
els dos nombres que formen el nombre combinatori i l'exponent de $B$ és el denominador del nombre combinatori. Amb aquestes observacions
previas no és difícil comprender que el terme que ocupal lloc $k$ ($1\leq k\leq n+1$) en el desenvolupament de $(A+B)^{n}$ sigui\begin{equation*}
t_{k}=\binom{n}{k-1}\cdot A^{n-(k-1)}\cdot B^{k-1}
\end{equation*}

Per exemple, es tracta de calcular el terme $5^{o}$ del desenvolupament de $(x^{2}-3x)^{6}$ sense necessitat de fer tot el desenvolupament. Així, tenim\begin{eqnarray*}
t_{5} &=&\binom{6}{4}\cdot \left( x^{2}\right) ^{6-4}\cdot \left( -3x\right)
^{4} \\
&=&15\cdot x^{4}\cdot 81x^{4} \\
&=&1215x^{8}
\end{eqnarray*}Observa que hem pres $A=x^{2}$ i $B=-3x$, perquè $(x^{2}-3x)^{6}=\left[
x^{2}+\left( -3x\right) \right] ^{6}$.

\section{Divisió de polinomis}

Donats dos polinomis qualssevol $p(x)$ i $q(x)$ tals que
grad\,$p(x)\geq $grad\,$q(x)$, es pot demostrar que sempre
existeixen dos polinomis $c(x)$ i $r(x)$ tals que\begin{equation*}
p(x)=q(x)\cdot c(x)+r(x)
\end{equation*}amb grad\,$r(x)<$ grad\,$q(x)$ o $r(x)=0$. L'operació que
consisteix en trobar els polinomis $c(x)$ i $r(x)$ s’anomena
\textbf{divisió} de $p(x)$ per $q(x)$. Aleshores els
polinomis $c(x)$ i $r(x)$ s'anomenen \textbf{quocient} i
\textbf{residu} de la divisió respectivament, i a $p(x)$ i
$q(x)$, \textbf{dividend} i \textbf{divisor} de la divisió,
respectivament.

Observa que, com que són grad\,$r(x)<$ grad\,$q(x)$, el grau del
quocient d’una divisió de polinomis sempre és la diferència
dels graus del dividend i divisor. En efecte, atès que
\begin{eqnarray*}
\text{grad\,}p(x) &=&\text{grad\,}\left( q(x)\cdot c(x)+r(x)\right) \\
&=&\text{grad\,}\left( q(x)\cdot c(x)\right) \text{ja que }\operatorname{grau}\,r(x)<\text{grad}\,q(x) \\
&=&\text{grad}\,q(x)+\text{grad}\,c(x)
\end{eqnarray*}d'on\begin{equation*}
\text{grad}\,c(x)=\text{grad}\,p(x)-\text{grad}\,q(x)
\end{equation*}

\subsection{Algorisme de la divisió}

Exposem a continuació un mètode per a trobar el quocient i el
residu d’una divisió de polinomis.

\begin{enumerate}
\item Ordenem els polinomis $p(x)$ i $q(x)$, dividend i divisor de la
divisió, en forma de potències decreixents en lloc de creixents com
fins a ara. El monomi de grau més gran del quocient s’obté dividint el
monomi de grau més gran de $p(x)$ pel monomi de grau més gran de $q(x)$.
Aleshores, es multiplica el divisor per aquest monomi i es resta aquest producte
del dividend. La diferència resultant és el primer residu parcial.

\begin{equation*}
\begin{array}{ccccccc|ccccc}
-2x^{3} & + & 3x^{2} & + & 0x & - & 1 & 5x^{2} & - & 3x & + & 2 \\
\cline{8-12}
2x^{3} & - & \frac{6}{5}x^{2} & + & \frac{4}{5}x & & & -\frac{2}{5}x & &
& & \\ \cline{1-5}
\diagup & & \frac{9}{5}x^{2} & + & \frac{4}{5}x & - & 1 & & & & &
\end{array}\end{equation*}

\item Si el residu parcial obtingut abans és de grau igual o més gran que el
del divisor, repetim el pas anterior fins a obtenir un residu de grau
més petit que el del divisor.

\begin{equation*}
\begin{array}{ccccccc|ccccc}
-2x^{3} & + & 3x^{2} & + & 0x & - & 1 & 5x^{2} & - & 3x & + & 2 \\
\cline{8-12}
2x^{3} & - & \frac{6}{5}x^{2} & + & \frac{4}{5}x & & & -\frac{2}{5}x & + &
\frac{9}{25} & & \\ \cline{1-5}
\diagup & & \frac{9}{5}x^{2} & + & \frac{4}{5}x & - & 1 & & & & & \\
& - & \frac{9}{5}x^{2} & + & \frac{27}{25}x & - & \frac{18}{25} & & & &
& \\ \cline{2-7}
& & \diagup & & \frac{47}{25}x & - & \frac{43}{25} & & & & &
\end{array}\end{equation*}Observa que el residu és ara\begin{equation*}
\frac{47}{25}x-\frac{43}{25}
\end{equation*}de grau més petit que el del divisor i, per tant, la divisió ha
finalizado.
\end{enumerate}

Com a conseqüència tenim que en dividir $p(x)=-2x^{3}+3x^{2}-1$ per $q(x)=5x^{2}-3x+2$ s'obtenen com residu i quocient, respectivament,\begin{equation*}
r(x)=\frac{47}{25}x-\frac{43}{25}\text{ \ \ \ i \ \ \ }c(x)=-\frac{2}{5}x+\frac{9}{25}
\end{equation*}i es compleix\begin{equation*}
-2x^{3}+3x^{2}-1=(5x^{2}-3x+2)\left( -\frac{2}{5}x+\frac{9}{25}\right) +\frac{47}{25}x-\frac{43}{25}
\end{equation*}

\subsection{Divisibilitat de polinomis}

Donats dos polinomis $p(x)$ i $q(x)$, diem que $p(x)$ és \textbf{divisible} per $q(x)$ o també que $q(x)$ és un \textbf{divisor} de $p(x)$ si
existeix un polinomi $c(x)$ tal que\begin{equation*}
p(x)=q(x)\cdot c(x)
\end{equation*}Aquesta relació també s'expressa dient que $p(x)$ és un \textbf{múltiple} de $q(x)$. Observa que si $p(x)$ és un múltiple de $q(x)$,
aleshores el residu de la divisió de $p(x)$ per $q(x)$ és exacta, és a
dir, es té residu igual a $0$.

Per exemple, $x^{2}-1$ és un múltiple de $x+1$, perquè\begin{equation*}
x^{2}-1=(x+1)\cdot (x-1)
\end{equation*}

Un polinomi de grau $n$
\begin{equation*}
p(x)=a_{0}+a_{1}x+a_{2}x^{2}+\cdots +a_{n-1}x^{n-1}+a_{n}x^{n}
\end{equation*}s’anomena \textbf{mònic} o \textbf{unitari} si el coeficient del
monomi de grau més gran és igual a la unitat, és a dir, $a_{n}=1$. Per exemple, els polinomis $x^{2}+1$, $2-3x+x^{2}$ o $1+x^{4}$, són mònics.

Tot polinomi pot escriure's com a producte d'un polinomi mònic del
mateix grau per un nombre real, perquè\begin{equation*}
a_{0}+a_{1}x+a_{2}x^{2}+\cdots +a_{n-1}x^{n-1}+a_{n}x^{n}=a_{n}\cdot \left(
\frac{a_{0}}{a_{n}}+\frac{a_{1}}{a_{n}}x+\frac{a_{2}}{a_{n}}x^{2}+\cdots +\frac{a_{n-1}}{a_{n}}x^{n-1}+x^{n}\right)
\end{equation*}
Per exemple, el polinomi $p(x)=1-3x-4x^{3}$ pot escriure's de la manera següent\begin{equation*}
p(x)=\left( -4\right) \cdot \left( -\frac{1}{4}+\frac{3}{4}x+x^{3}\right)
\end{equation*}i, en aquest cas, es diu que el polinomi\begin{equation*}
-\frac{1}{4}+\frac{3}{4}x+x^{3}
\end{equation*}és el polinomi mònic de $p(x)$.

Un polinomi $p(x)$ s’anomena \textbf{irreduïble} si els únics
divisors mònics de $p(x)$ són el polinomi unitat $1$ i el seu polinomi mònic corresponent. Per exemple, els polinomis $x+1$, $x^{2}+1$, $2x-3$ o $x^{2}+x+1$ són irreduïbles.

\subsection{Màxim comú divisor de dos polinomis}

S'anomena \textbf{màxim comú divisor} dels polinomis $p(x)$ i $q(x)$ el polinomi mònic $d(x)$ que satisfà la propietat següent: $d(x)$ és el polinomi de grau més gran que és divisor de $p(x)$ i de $q(x)$.
Escriurem\begin{equation*}
d(x)=\text{m.c.d.}\left( p(x),q(x)\right)
\end{equation*}

Per exemple, el màxim comú divisor dels polinomis $p(x)=x^{2}+2x+1$ i $q(x)=x^{2}-1$ és $d(x)=x+1$, perquè,\begin{eqnarray*}
p(x) &=&(x+1)^{2} \\
q(x) &=&(x+1)(x-1)
\end{eqnarray*}i, per tant, $x+1$ el divisor comú de grau més gran. Així,
escriurem\begin{equation*}
\text{m.c.d.}\left( x^{2}+2x+1,x^{2}-1\right) =x+1
\end{equation*}

Si $p(x)$ i $q(x)$ són dos polinomis tals que $\operatorname{grau}~p(x)\geq $$\operatorname{grau}~q(x)$, aleshores es compleix\begin{equation*}
\text{m.c.d.}\left( p(x),q(x)\right) =\text{m.c.d.}\left( q(x),r(x)\right)
\end{equation*}on $r(x)$ és el residu de la divisió de $p(x)$ per $q(x)$.

En efecte, dividint $p(x)$ per $q(x)$ obtenim\begin{equation}
p(x)=q(x)\cdot c(x)+r(x) \label{1}
\end{equation}i, per tant,\begin{equation}
r(x)=p(x)-q(x)\cdot c(x) \label{2}
\end{equation}Com que $d(x)=$~m.c.d.$\left( p(x),q(x)\right) $ és un divisor de $p(x)$
i de $q(x)$, aleshores, per (\ref{2}), $d(x)$ serà també un divisor
de $r(x)$. Com a conseqüència, $d(x)$ és un divisor de m.c.d.$\left(
q(x),r(x)\right) $. D'altra banda, i de la mateixa manera, com que m.c.d.$\left(
q(x),r(x)\right) $ és un divisor de $q(x)$ i $r(x)$, aleshores, per (\ref{1}), m.c.d.$\left( q(x),r(x)\right) $serà també un divisor de $p(x)$. Com a conseqüència, m.c.d.$\left( q(x),r(x)\right) $ és un divisor de $d(x)$. Hem demostrat així que $d(x)$ divideix a m.c.d.$\left(
q(x),r(x)\right) $i que m.c.d.$\left( q(x),r(x)\right) $divideix a $d(x)$.
Per tant, com que són tots dos polinomis mònics, es dedueix\footnote{És fàcil demostrar que donats dos polinomis mònics, si el primer
divideix al segon i el segon el primer, aleshores els dos polinomis són
iguals.}\begin{equation*}
\text{m.c.d.}\left( p(x),q(x)\right) =d(x)=\text{m.c.d.}\left(
q(x),r(x)\right)
\end{equation*}

Aquest resultat és la base de l'algorisme d'Euclides per a trobar el màxim comú divisor de dos polinomis.

\subsubsection{algorisme d’Euclides}

Donats dos polinomis $p(x)$ i $q(x)$ tals que $\operatorname{grau}~p(x)\geq $$\operatorname{grau}~q(x)$,
efectuem la divisió de $p(x)$ per $q(x)$ i obtenim un residu $r_{1}(x) $. Aleshores es compleix\begin{equation*}
\text{m.c.d.}\left( p(x),q(x)\right) =\text{m.c.d.}\left(
q(x),r_{1}(x)\right)
\end{equation*}amb $\operatorname{grau}~r_{1}(x)<$ $\operatorname{grau}~q(x)$. Si $r_{1}(x)$ no és un nombre real,
efectuem la divisió de $q(x)$ per $r_{1}(x)$ i obtenim un residu $r_{2}(x)$. Aleshores es compleix\begin{equation*}
\text{m.c.d.}\left( q(x),r_{1}(x)\right) =\text{m.c.d.}\left(
r_{1}(x),r_{2}(x)\right)
\end{equation*}amb $\operatorname{grau}~r_{2}(x)<$ $\operatorname{grau}~r_{1}(x)$. Si encara $r_{2}(x)$ no és un nombre real, continuem dividint fins a obtenir residu $r_{n}(x)$ que
sigui un nombre real; observa que això passarà amb seguretat perquè\begin{equation*}
\text{grau}~p(x)\geq \text{grau}~q(x)>\text{grau}~r_{1}(x)>\text{grau}~r_{2}(x)>\cdots
\end{equation*}Com a conseqüència, obtindrem\begin{eqnarray*}
\text{m.c.d.}\left( p(x),q(x)\right) &=&\text{m.c.d.}\left(
q(x),r_{1}(x)\right) \\
&=&\text{m.c.d.}\left( r_{1}(x),r_{2}(x)\right) \\
&&\cdots \\
&=&\text{m.c.d.}\left( r_{n-1}(x),r_{n}(x)\right) \\
&=&r_{n-1}^{\ast }(x)
\end{eqnarray*}essent $r_{n}(x)$ un nombre real i $r_{n-1}^{\ast }(x)$ el polinomi mònic de $r_{n-1}(x)$.

Per a efectuar les successives divisions i trobar $r_{n-1}(x)$ se sol
utilitzar el esquema següent\begin{equation*}
\begin{array}{c|c|c|c|c|c}
& c_{1}(x) & c_{2}(x) & & c_{n-1}(x) & \\ \hline
p(x) & q(x) & r_{1}(x) & \cdots & r_{n-2}(x) & r_{n-1}(x) \\ \hline
r_{1}(x) & r_{2}(x) & r_{3}(x) & & r_{n}(x)\in \mathbb{R} &
\end{array}\end{equation*}Per exemple, calcularem el màxim comú divisor dels polinomis\begin{eqnarray*}
p(x) &=&x^{5}+2x^{3}+3x^{2}-15x-9 \\
q(x) &=&x^{4}-9
\end{eqnarray*}segons el esquema considerat:\begin{equation*}
\begin{array}{ccccc|ccccc|cccc|cc}
& & & & & x & & & & & \frac{1}{2}x & -\frac{3}{4} & & & 2x & +3 \\
\cline{6-16}
x^{5} & +2x^{3} & +3x^{2} & -15x & -9 & x^{4} & 0 & 0 & 0 & -9 & 2x^{3} &
+3x^{2} & -6x & -9 & x^{2} & -3 \\
-x^{5} & & & +9x & & -x^{4} & -\frac{3}{2}x^{3} & +3x^{2} & +\frac{9}{2}x
& & -2x^{3} & & +6x & & & \\ \cline{1-5}\cline{6-9}\cline{11-13}
\diagup & +2x^{3} & +3x^{2} & -6x & -9 & \diagup & -\frac{3}{2}x^{3} &
+3x^{2} & +\frac{9}{2}x & -9 & \diagup & 3x^{2} & \diagup & -9 & & \\
& & & & & & \frac{3}{2}x^{3} & +\frac{9}{4}x^{2} & -\frac{9}{2}x & -\frac{27}{4} & & -3x^{2} & & +9 & & \\ \cline{7-10}\cline{12-14}
& & & & & & \diagup & \frac{21}{4}x^{2} & \diagup & -\frac{63}{4} & &
\diagup & & \diagup & & \\
& & & & & & & x^{2} & & -3 & & & & & &
\end{array}\end{equation*}és a dir, hem obtingut que\begin{equation*}
\text{m.c.d.}\left( x^{5}+2x^{3}+3x^{2}-15x-9,x^{4}-9\right) =\text{m.c.d.}\left( x^{2}-3,0\right) =x^{2}-3
\end{equation*}Observa que hem aplicat en el esquemal següent fet\begin{equation*}
\text{m.c.d.}\left( 2x^{3}+3x^{2}-6x-9,\frac{21}{4}(x^{2}-3)\right) =\text{m.c.d.}\left( 2x^{3}+3x^{2}-6x-9,x^{2}-3\right)
\end{equation*}cosa que és evident.

\subsection{Mínim comú múltiple de dos polinomis}

S'anomena \textbf{mínim comú múltiple} dels polinomis $p(x)$
i $q(x)$ el polinomi mònic $m(x)$ que satisfà la següent
propietat: $m(x)$ és el polinomi de grau més petit que és múltiple de $p(x) $i de $q(x)$. Escriurem\begin{equation*}
m(x)=\text{m.c.m.}\left( p(x),q(x)\right)
\end{equation*}

Per exemple, el mínim comú múltiple dels polinomis $p(x)=x^{2}+2x+1$ i $q(x)=x^{2}-1$ és $m(x)=(x+1)^{2}(x-1)$, perquè,\begin{eqnarray*}
p(x) &=&(x+1)^{2} \\
q(x) &=&(x+1)(x-1)
\end{eqnarray*}i, per tant, $(x+1)^{2}(x-1)$ el múltiple comú de més petit
grau. Així, escriurem\begin{equation*}
\text{m.c.m.}\left( x^{2}+2x+1,x^{2}-1\right) =(x+1)^{2}(x-1)
\end{equation*}

És clar que si $p(x)$ és múltiple de $q(x)$, aleshores el mínim comú múltiple és $p(x)$. Es demostra en cursos superiors que es
compleix la següent relació\begin{equation*}
\text{m.c.d.}\left( p(x),q(x)\right) \cdot \text{m.c.m.}\left(
p(x),q(x)\right) =\frac{p(x)\cdot q(x)}{a_{n}\cdot b_{m}}
\end{equation*}on $a_{n}$ i $b_{m}$ són els coeficients dels monomis de grau més gran
de $p(x)$ i~$q(x)$, respectivament. Aquesta relació nos sirve per a trobar
el mínim comú múltiple a partir del màxim comú
divisor, perquè,\begin{equation*}
\text{m.c.m.}\left( p(x),q(x)\right) =\frac{1}{a_{n}\cdot b_{m}}\frac{p(x)\cdot q(x)}{\text{m.c.d.}\left( p(x),q(x)\right) }
\end{equation*}Per exemple, si
\begin{eqnarray*}
p(x) &=&x^{5}+2x^{3}+3x^{2}-15x-9 \\
q(x) &=&x^{4}-9
\end{eqnarray*}aleshores\begin{eqnarray*}
\text{m.c.m.}\left( x^{5}+2x^{3}+3x^{2}-15x-9,x^{4}-9\right) &=&\frac{(x^{5}+2x^{3}+3x^{2}-15x-9)(x^{4}-9)}{x^{2}-3} \\
&=&x^{7}+5x^{5}+3x^{4}-9x^{3}-45x-27
\end{eqnarray*}

En un altre apartat veurem com la descomposició factorial d'un
polinomi dona un altre mètode de càlcul del màxim comú divisor
i mínim comú múltiple.

\section{Regla de Ruffini i teorema del residu}

En aquest apartat veurem com es poden calcular el quocient i el residu de
la divisió d'un polinomi per un binomi de primer grau de la forma $x-a$ per un mètode més senzill que l'aplicació de l'algorisme
clàssic de la divisió. Aquest mètode es coneix com la regla de
Ruffini. També donarem un teorema amb el que podrem saber si un
polinomi és o no divisible per un binomi de primer grau de la forma $x-a$
amb només calcular el valor numèric del polinomi per a $x=a$ i
veure si és o no zero. Aquest resultat es coneix com el teorema del residu.

\subsection{Regla de Ruffini}

Donat un polinomi qualsevol $p(x)$, ara volem dividir-lo pel binomi
de primer grau $x-a$, és a dir, volem trobar els polinomis quocient $c(x)$ i residu $r(x)$ tals que\begin{equation}
p(x)=(x-a)\cdot c(x)+r(x) \label{3}
\end{equation}Observa primer que $\operatorname{grau}~r(x)=0$ o bé $r(x)=0$, perquè, sabem que $\operatorname{grau}~r(x)<$ $\operatorname{grau}~(x-a)=1$. Per tant, $r(x)$ serà un nombre real $r$.
Observa també que\begin{equation*}
\text{grau}~c(x)=\text{grau}~p(x)-\text{grau}~(x-a)=\text{grau}~p(x)-1
\end{equation*}Per tant, si\begin{equation*}
p(x)=a_{0}+a_{1}x+a_{2}x^{2}+\cdots +a_{n-1}x^{n-1}+a_{n}x^{n}
\end{equation*}aleshores\begin{equation*}
c(x)=b_{0}+b_{1}x+b_{2}x^{2}+\cdots +b_{n-2}x^{n-2}+b_{n-1}x^{n-1}
\end{equation*}i segons (\ref{3}), tenim\begin{eqnarray*}
a_{0}+a_{1}x+a_{2}x^{2}+\cdots +a_{n-1}x^{n-1}+a_{n}x^{n}
&=&(x-a)(b_{0}+b_{1}x+b_{2}x^{2}+\cdots +b_{n-2}x^{n-2}+b_{n-1}x^{n-1})+r \\
&=&(r-ab_{0})+(b_{0}-ab_{1})x+(b_{1}-ab_{2})x^{2}+\cdots + \\
&&(b_{n-2}-ab_{n-1})x^{n-1}+b_{n-1}x^{n}
\end{eqnarray*}i, per tant,\begin{equation*}
\begin{array}{l}
a_{0}=r-ab_{0} \\
a_{1}=b_{0}-ab_{1} \\
a_{2}=b_{1}-ab_{2} \\
\multicolumn{1}{c}{\cdots} \\
a_{n-1}=b_{n-2}-ab_{n-1} \\
a_{n}=b_{n-1}\end{array}\end{equation*}i, en conseqüència, obtenim les següents fórmules que ens permetran determinar $c(x)$ i $r$\begin{equation*}
\begin{array}{l}
b_{n-1}=a_{n} \\
b_{n-2}=a_{n-1}+a\cdot b_{n-1} \\
\multicolumn{1}{c}{\cdots} \\
b_{1}=a_{2}+ab_{2} \\
b_{0}=a_{1}+ab_{1} \\
r=a_{0}+ab_{0}\end{array}\end{equation*}Tanmateix, el càlcul dels coeficients $b_{i}$ ($0\leq i\leq n-1$)
i $r$ es poden fer de manera molt senzilla mitjançant l'esquema següent
anomenat \textbf{regla de Ruffini}.

\begin{center}
\includegraphics[width=4.0906in,height=0.9331in,keepaspectratio]{poli1.png}
\end{center}

Per exemple, calcularem el quocient i el residu de dividir $p(x)=x^{3}-2x+7$
per $x-2$. Per fer-ho, farem servir la regla de Ruffini.\begin{equation*}
\begin{array}{c|cccc}
& 1 & 0 & -2 & 7 \\
2 & & 2 & 4 & 4 \\ \hline
& 1 & 2 & 2 & 11\end{array}\end{equation*}Dit amb paraules, baixem el primer coeficient del dividend. Aleshores,
multipliquem per 2 i col·loquem el resultat sota del segon coeficient
del dividend i, després, els sumem. Repetim aquest procés fins a
completar el quadre. L'última fila del quadre són els coeficients del
quocient (el seu grau és una unitat inferior al del dividend) i l'últim nombre és el residu. Així, tenim $c(x)=x^{2}+2x+2$ i $r=11$.

Observa que es col·loquen zeros per als coeficients dels monomis que no
apareixen en el polinomi dividend.

\subsection{Teorema del residu}

El valor numèric del polinomi $p(x)$ per a $x=a$ coincideix amb el residu
de dividir $p(x)$ per $x-a$.

En efecte, si dividim $p(x)$ per $x-a$ obtenim un quocient $c(x)$ i un
residu $r$ de forma que\begin{equation*}
p(x)=(x-a)\cdot c(x)+r
\end{equation*}Aleshores, el valor numèric de $p(x)$ per a $x=a$ és\begin{equation*}
p(a)=(a-a)\cdot c(a)+r=0+r=r
\end{equation*}és a dir, es compleix que $p(a)=r$.

Per exemple, calcularem el residu $r$ de dividir $x^{3}-2x+7$ per $x-2$. Segons el teorema del residu es compleix\begin{equation*}
r=p(2)=2^{3}-2\cdot 2+7=11
\end{equation*}

Aleshores, observa que un polinomi serà divisible pel binomi de
primer grau $x-a$ si el valor numèric del polinomi és zero per a $x=a$.

Per exemple, volem esbrinar si el polinomi $p(x)=x^{3}-2x+1$ es
divisible per $x-1$. Per a contestar a aquesta pregunta només cal
calcular el valor numèric de $p(x)$ per a $x=1$ i veure si és zero o no.
En concret,\begin{equation*}
p(1)=1-2+1=0
\end{equation*}Per tant, $p(x)$ és divisible per \thinspace $x-1$.

\section{Arrels i descomposició factorial d'un polinomi}

Diem que un nombre real $a$ és una \textbf{arrel} o \textbf{zero}
d'un polinomi $p(x)$ si $p(a)=0$. Així, $1$ és una arrel del
polinomi $p(x)=x^{2}-1$ perquè $p(1)=0$, així com $\sqrt{3}$ és una arrel de $q(x)=x^{2}-3$, perquè,\begin{equation*}
q(\sqrt{3})=\left( \sqrt{3}\right) ^{2}-3=0
\end{equation*}

Segons el teorema del residu, si $a$ és una arrel de $p(x)$, aleshores $p(x)$ és divisible per $x-a$ i, per tant, podem escriure\begin{equation*}
p(x)=(x-a)\cdot c(x)
\end{equation*}

Aleshores, diem que $a$ és una \textbf{arrel simple} de $p(x)$ si $c(a)\neq 0$, és a dir, quan el quocient de dividir $p(x)$ per $x-a$ no
torna a ser divisible per $x-a$.

Per exemple, $3$ és una arrel simple del polinomi $p(x)=x^{3}+x^{2}-9x-9 $, ja que per Ruffini\begin{equation*}
\begin{array}{r|rrrr}
& 1 & 1 & -9 & -9 \\
3 & & 3 & 12 & 9 \\ \hline
& 1 & 4 & 3 & 0\end{array}\end{equation*}obtenim\begin{equation*}
p(x)=(x-3)(x^{2}+4x+3)
\end{equation*}i a més es compleix\begin{equation*}
c(3)=3^{2}+4\cdot 3+3=24\neq 0
\end{equation*}

Diem que $a$ és una \textbf{arrel múltiple} de $p(x)$ si $a$ es
una arrel simple
\begin{equation*}
p(x)=(x-a)\cdot c(x)
\end{equation*}i a més es compleix $c(a)=0$, és a dir, si $c(x)$ torna a ser divisible
per $x-a$.

Per exemple, $1$ és una arrel múltiple del polinomi $p(x)=x^{3}+x^{2}-5x+3$, perquè, per Ruffini\begin{equation*}
\begin{array}{r|rrrr}
& 1 & 1 & -5 & 3 \\
1 & & 1 & 2 & -3 \\ \hline
& 1 & 2 & -3 & 0\end{array}\end{equation*}obtenim\begin{equation*}
p(x)=(x-1)(x^{2}+2x-3)
\end{equation*}i es compleix\begin{equation*}
c(1)=1+2-3=0
\end{equation*}A més, observa que si ara tornem a aplicar Ruffini al quocient,\begin{equation*}
\begin{array}{r|rrr}
& 1 & 2 & -3 \\
1 & & 1 & 3 \\ \hline
& 1 & 3 & 0\end{array}\end{equation*}aleshores obtenim\begin{equation*}
c(x)=(x-1)(x+3)
\end{equation*}i, com a conseqüència,\begin{eqnarray*}
p(x) &=&(x-1)(x-1)(x+3) \\
&=&(x-1)^{2}(x+3)
\end{eqnarray*}Observem que en aquest cas hem pogut dividir dos vegades $p(x)$ per $x-1$. Aleshores, diem que $1$ és una arrel múltiple de ordre 2 o
simplement una arrel doble.

En general, si $a$ és una arrel de $p(x)$ i es compleix\begin{equation*}
p(x)=(x-a)^{s}\cdot c(x)
\end{equation*}amb $c(a)\neq 0$, aleshores diem que $a$ és una \textbf{arrel múltiple de ordre} $s$ o també que $a$ és una \textbf{arrel de
multiplicitat} $s$.

Per exemple, volem esbrinar la multiplicitat de l'arrel $-1$ en el
polinomi\begin{equation*}
p(x)=x^{5}-x^{4}-5x^{3}+x^{2}+8x+4
\end{equation*}Per fer-ho, aplicarem repetidament la regla de Ruffini\begin{equation*}
\begin{array}{r|rrrrrr}
& 1 & -1 & -5 & 1 & 8 & 4 \\
-1 & & -1 & 2 & 3 & -4 & -4 \\ \hline
& 1 & -2 & -3 & 4 & 4 & 0 \\
-1 & & -1 & 3 & 0 & -4 & \\ \hline
& 1 & -3 & 0 & 4 & 0 & \\
-1 & & -1 & 4 & -4 & & \\ \hline
& 1 & -4 & 4 & 0 & & \\
-1 & & -1 & 5 & & & \\ \hline
& 1 & -5 & 9 & & &
\end{array}\end{equation*}
Observem que $p(x)$ pot dividir-se 3 vegades per $x+1$. A més, es
compleix\begin{equation*}
p(x)=(x+1)^{3}(x^{2}-4x+4)
\end{equation*}i el tercer quocient que s’obté de la tercera divisió per $x+1$
compleix\begin{equation*}
c_{3}(-1)=1+4+4=9\neq 0
\end{equation*}Per tant, $-1$ és una arrel de $p(x)$ de multiplicitat $3$.

\subsection{Arrels enteres d'un polinomi amb coeficients enters}

Considerem ara un polinomi $p(x)$ amb coeficients enters, és a dir,\begin{equation*}
p(x)=a_{0}+a_{1}x+a_{2}x^{2}+\cdots +a_{n-1}x^{n-1}+a_{n}x^{n}
\end{equation*}i $a_{i}\in \mathbb{Z}$ ($0\leq i\leq n$). Suposem que $k$ és una arrel entera ($k\in \mathbb{Z}$) de $p(x)$. Aleshores, es compleix\begin{eqnarray*}
a_{0}+a_{1}k+a_{2}k^{2}+\cdots +a_{n-1}k^{n-1}+a_{n}k^{n} &=&0 \\
(a_{1}+a_{2}k+\cdots +a_{n-1}k^{n-2}+a_{n}k^{n-1})~k &=&-a_{0} \\
a_{1}+a_{2}k+\cdots +a_{n-1}k^{n-2}+a_{n}k^{n-1} &=&-\frac{a_{0}}{k}
\end{eqnarray*}i com $a_{1}+a_{2}k+\cdots +a_{n-1}k^{n-2}+a_{n}k^{n-1}\in \mathbb{Z}$, $k$
ha de ser necessàriament un divisor del terme independent $a_{0}$.
D'aquesta manera, les arrels enteres d'un polinomi amb coeficients
enters s'han de trobar entre els divisors del seu terme
independent.

Per exemple, si volem calcular les arrels enteres del polinomi $p(x)=x^{4}-2x^{3}-2x^{2}-2x+3$, caldrà provar només amb $\pm 1$ i $\pm
3$, que són els divisors del terme independent. Un càlcul senzill mostra que $3$ i $-1$ són les úniques arrels enteres d'aquest
polinomi.

\subsection{Descomposició factorial d'un polinomi}

De la mateixa manera que tot nombre enter pot descompondre's com a producte
de nombres primers, tot polinomi pot descompondre's en producte de
polinomis irreduïbles. La descomposició varia segons el
anell de polinomis que considerem. Per exemple, $x^{2}-2$ és irreduïble
en $\mathbb{Q}[x]$ i, en canvi, pot descompondre's en
\begin{equation*}
(x-\sqrt{2})(x+\sqrt{2})
\end{equation*}
en $\mathbb{R}[x]$. Admetem, sense demostrar-ho, que els polinomis
irreduïbles en $\mathbb{R}[x]$ són de primer grau o bé de segon grau.
Qualsevol polinomi de primer grau $a+bx$ és irreduïble, perquè no té
altres divisors mònics més que la unitat i el polinomi\begin{equation*}
\frac{a}{b}+x
\end{equation*}Un polinomi de segon grau és irreduïble si l'equació $c+bx+ax^{2}=0$ no té arrels reals, és a dir, quan $b^{2}-4ac<0$.

Ja hem vist en l'apartat anterior que si $k_{1}$ és una arrel de $p(x)$, aleshores\begin{equation*}
p(x)=(x-k_{1})\cdot c_{1}(x)
\end{equation*}Si ara $c_{1}(x)$ admet una altra arrel $k_{2}$, també serà arrel de $p(x)$ i es compleix\begin{equation*}
p(x)=(x-k_{1})\cdot (x-k_{2})\cdot c_{2}(x)
\end{equation*}Reiterant aquest procés, si $p(x)$ admet $n$ arrels reals $k_{1},k_{2},...,k_{n}$, aleshores podrem descompondre $p(x)$ en factors
lineals (polinomis de primer grau) de la forma següent\begin{equation}
p(x)=a_{n}\cdot (x-k_{1})\cdot (x-k_{2})\cdot \cdots \cdot (x-k_{n})
\label{4}
\end{equation}on $a_{n}$ és el coeficient del monomi de grau $n$ de $p(x)$.

Observa que un polinomi de grau $n$ no pot tenir més que $n$ arrels reals diferents, perquè, si no fos així, si $k_{n+1}$ fos una altra arrel real diferent a totes les anteriors, aleshores caldria complir-se\begin{equation*}
p(k_{n+1})=a_{n}\cdot (k_{n+1}-k_{1})\cdot (k_{n+1}-k_{2})\cdot \cdots \cdot
(k_{n+1}-k_{n})=0
\end{equation*}cosa que implicaria $a_{n}=0$, i això és contradictori amb el fet de
que el polinomi és de grau $n$.

Observa que no sempre podrà efectuar-se una descomposició com (\ref{4}), perquè, com ja hem dit, existeixen polinomis de segon grau que són
irreduïbles i, per tant, no poden descompondre's en factors lineals. Per exemple,
\begin{equation*}
x^{3}-1=(x+1)\cdot (x^{2}-x+1)
\end{equation*}i el factor $x^{2}-x+1$ no pot descompondre's més ja que $b^{2}-4ac=-3<0 $.

En general, podrem obtenir una descomposició d'un polinomi $p(x)$
de la forma\begin{equation}
p(x)=p_{1}(x)\cdot p_{2}(x)\cdot \cdots \cdot p_{r}(x) \label{5}
\end{equation}amb\begin{equation*}
\text{grau}~p(x)=\text{grau}~p_{1}(x)+\text{grau}~p_{2}(x)+\cdots +\text{grau}~p_{r}(x)
\end{equation*}on cada $p_{i}(x)$ ($1\leq i\leq r$) és un polinomi irreduïble de $\mathbb{R}[x]$.

Per exemple, el polinomi $p(x)=x^{4}+x^{3}+2x^{2}+x+1$ admet la següent
descomposició
\begin{equation*}
p(x)=(x^{2}+1)\cdot (x^{2}+x+1)
\end{equation*}els factors del qual són polinomis irreduïbles.

Si cada $p_{i}(x)$ és lineal, és a dir, si $p(x)$ admet totes les arrels i són simples, aleshores (\ref{5}) s'escriurà com\begin{equation*}
p(x)=a_{n}\cdot (x-k_{1})\cdot (x-k_{2})\cdot \cdots \cdot (x-k_{n})
\end{equation*}essent $k_{i}$ ($1\leq i\leq n$) les arrels de $p(x)$.

Per exemple, és immediat comprovar que el polinomi $2x^{4}-10x^{2}+8$
té arrels $\pm 1$ i $\pm 2$ i, per tant, es descompon en factors
lineals de la forma següent\begin{equation*}
2x^{4}-10x^{2}+8=2\cdot (x-1)\cdot (x+1)\cdot (x-2)\cdot (x+2)
\end{equation*}

Si el polinomi $p(x)$ admet totes les arrels i algunes són múltiples, agrupant els factors lineals iguals, obtindrem una
descomposició de la forma\begin{equation*}
p(x)=a_{n}\cdot (x-k_{1})^{n_{1}}\cdot (x-k_{2})^{n_{2}}\cdot \cdots \cdot
(x-k_{r})^{n_{r}}
\end{equation*}amb $n_{1}+n_{2}+\cdots +n_{r}=n$, essent $n_{i}$ ($1\leq i\leq r$) la
multiplicitat de l'arrel $k_{i}$.

Per exemple, per a descompondre el polinomi $p(x)=x^{5}-5x^{4}+3x^{3}+13x^{2}-8x-12$ podem utilitzar la regla de
Ruffini. A més, com els seus coeficients són enters, les arrels seran divisors del terme independent, és a dir, es trobaran entre $\pm 1$, $\pm 2$, $\pm 3$, $\pm 4$, $\pm 6$ i $\pm 12$. Observa que es compleix\begin{equation*}
\begin{array}{r|rrrrrr}
& 1 & -5 & 3 & 13 & -8 & -12 \\
-1 & & -1 & 6 & -9 & -4 & 12 \\ \hline
& 1 & -6 & 9 & 4 & -12 & 0 \\
-1 & & -1 & 7 & -16 & 12 & \\ \hline
& 1 & -7 & 16 & -12 & 0 & \\
2 & & 2 & -10 & 12 & & \\ \hline
& 1 & -5 & 6 & 0 & & \\
2 & & 2 & -6 & & & \\ \hline
& 1 & -3 & 0 & & & \\
3 & & 3 & & & & \\ \hline
& 1 & 0 & & & &
\end{array}\end{equation*}o sigui, les arrels de $p(x)$ són $-1$ (doble), $2$ (doble) i $3$. Per tant, $p(x)$ es descompondrà de la forma següent\begin{equation*}
p(x)=(x+1)^{2}(x-2)^{2}(x-3)
\end{equation*}

\subsubsection{Càlcul del m.c.d. i m.c.m.}

Conocidas les descomposicions factorials en $\mathbb{R}[x]$ de dos
polinomis $p(x)$ i $q(x)$, el màxim comú divisor s'obtindrà
multiplicant els factors irreduïbles comuns elevats a l'exponent més petit que figuri en les dues descomposicions, i el mínim comú múltiple s'obtindrà multiplicant els factors irreduïbles
comuns i no comuns elevats a l'exponent més gran que figuri en
totes dues descomposicions.

Per exemple, calcularem el m.c.d. i el m.c.m. dels polinomis\begin{eqnarray*}
p(x) &=&2x^{5}-2x^{4}-10x^{3}+2x^{2}+16x+8 \\
q(x) &=&3x^{3}-6x^{2}-3x+6
\end{eqnarray*}Primer, buscarem les descomposicions de tots dos polinomis. Els polinomis
mònics de cada polinomi són\begin{eqnarray*}
p^{\ast }(x) &=&x^{5}-x^{4}-5x^{3}+x^{2}+8x+4 \\
q^{\ast }(x) &=&x^{3}-2x^{2}-x+2
\end{eqnarray*}
Aleshores, per a $p^{\ast }(x)$ tenim\begin{equation*}
\begin{array}{r|rrrrrr}
& 1 & -1 & -5 & 1 & 8 & 4 \\
-1 & & -1 & 2 & 3 & -4 & -4 \\ \hline
& 1 & -2 & -3 & 4 & 4 & 0 \\
-1 & & -1 & 3 & 0 & -4 & \\ \hline
& 1 & -3 & 0 & 4 & 0 & \\
-1 & & -1 & 4 & -4 & & \\ \hline
& 1 & -4 & 4 & 0 & & \\
2 & & 2 & -4 & & & \\ \hline
& 1 & -2 & 0 & & &
\end{array}\end{equation*}i, per tant,\begin{equation*}
p^{\ast }(x)=(x+1)^{3}\cdot (x-2)^{2}
\end{equation*}i\begin{equation*}
p(x)=2\cdot (x+1)^{3}\cdot (x-2)^{2}
\end{equation*}Per a $q^{\ast }(x)$ tenim\begin{equation*}
\begin{array}{r|rrrr}
& 1 & -2 & -1 & 2 \\
-1 & & -1 & 3 & -2 \\ \hline
& 1 & -3 & 2 & 0 \\
1 & & 1 & -2 & \\ \hline
& 1 & -2 & 0 &
\end{array}\end{equation*}i, per tant,\begin{equation*}
q^{\ast }(x)=(x+1)\cdot (x-1)\cdot (x-2)
\end{equation*}i\begin{equation*}
q(x)=3\cdot (x+1)\cdot (x-1)\cdot (x-2)
\end{equation*}Ara ja podem determinar m.c.d. i m.c.m. de $p(x)$ i $q(x)$. Així,
tenim\begin{equation*}
\text{m.c.d.}\left( p(x),q(x)\right) =(x+1)\cdot (x-2)
\end{equation*}i\begin{equation*}
\text{m.c.m.}\left( p(x),q(x)\right) =(x+1)\cdot (x+1)^{3}\cdot (x-2)^{2}
\end{equation*}

\section{Fraccions racionals}

De la mateixa manera que l'equació
\begin{equation*}
3\cdot x=2
\end{equation*}no té solució en l'anell dels nombres enters i per a
resoldre aquest problema tuvimos que ampliar $\mathbb{Z}$ construint el
conjunt dels nombres racionals, l'equació\begin{equation*}
(x-1)\cdot A(x)=x^{2}+1
\end{equation*}no té solució en l'anell de polinomis $\mathbb{R}[x]$ (ja que $x^{2}+1$ no és múltiple de $x-1$) i per a solucióar aquest problema
caldrà ampliar $\mathbb{R}[x]$ construint el conjunt de les fraccions
racionals.

S'anomena \textbf{fracció algebraica} o \textbf{polinòmica} a tota
expressió de la forma\begin{equation*}
\frac{p(x)}{q(x)}
\end{equation*}on $p(x)$ i $q(x)$ són polinomis i $q(x)$ és diferent del polinomi zero.

Per exemple,\begin{equation*}
\frac{x^{2}+1}{x},\frac{x-1}{2},\frac{x^{3}}{x^{2}-1}
\end{equation*}són fraccions algebraiques.

En el conjunt de totes les fraccions algebraiques definim una relació d'equivalència de la forma següent: diem que dues fraccions
algebraiques\begin{equation*}
\frac{p(x)}{q(x)}\text{ \ \ \ i \ \ \ }\frac{r(x)}{s(x)}
\end{equation*}són \textbf{equivalents} si i només si\begin{equation*}
p(x)\cdot s(x)=q(x)\cdot r(x)
\end{equation*}

Aquesta relació d'equivalència dona lloc a una clesificación de les
fraccions algebraiques en classes. Cada fracció
\begin{equation*}
\frac{p(x)}{q(x)}
\end{equation*}define una classe, que es pot representar per\begin{equation*}
\left\{ \frac{p(x)}{q(x)}\right\}
\end{equation*}a la qual pertanyen totes les fraccions equivalents, és a dir,\begin{equation*}
\left\{ \frac{p(x)}{q(x)}\right\} =\left\{ \frac{r(x)}{s(x)}:p(x)\cdot
s(x)=q(x)\cdot r(x)\right\}
\end{equation*}Anomenem \textbf{fracció racional} cada classe de fraccions
algebraiques. Com resulta molt feixuc escriure cada fracció racional
mitjançant claus, convindrem a escriurela simplement per un representant
de la classe. Així, en lloc de escriure\begin{equation*}
\left\{ \frac{p(x)}{q(x)}\right\}
\end{equation*}escriurem\begin{equation*}
\frac{p(x)}{q(x)}
\end{equation*}prenent aquesta fracció com a representant de la fracció racional. De
aquest modo, si\begin{equation*}
\frac{p(x)}{q(x)}\text{ \ \ \ i \ \ \ }\frac{r(x)}{s(x)}
\end{equation*}són equivalents, atès que pertanyen a la mateixa classe i cada una delles
representa una fracció racional, escriurem\begin{equation*}
\frac{p(x)}{q(x)}=\frac{r(x)}{s(x)}
\end{equation*}Així, amb aquest conveni,
\begin{equation*}
\frac{x-1}{x^{2}}=\frac{x^{2}-1}{x^{2}(x-1)}
\end{equation*}no significa que siguin iguals les fraccions algebraiques, sinó que són
iguals les fraccions racionals que determinen.

Una fracció algebraica
\begin{equation*}
\frac{p(x)}{q(x)}
\end{equation*}s’anomena \textbf{irreduïble} si $p(x)$ i $q(x)$ són primers entre si,
és a dir, si no té cap factor comú de grau més gran o igual
que $1$. Per exemple, les fraccions següents són irreduïbles\begin{equation*}
\frac{x-1}{x^{2}},\frac{x^{2}+1}{x-1},\frac{x^{3}}{x^{2}+1}
\end{equation*}mentre que les fraccions següents\begin{equation*}
\frac{x}{x^{2}-x},\frac{x^{2}-1}{x+1},\frac{x^{3}+1}{x^{2}-1}
\end{equation*}no ho són.

S'enté per \textbf{simplificar} una fracció algebraica buscar
la fracció irreduïble equivalent. En la pràctica i per a facilitar
els càlculs, sovint es pren com a representant d’una fracció
racional la fracció irreduïble.

Per exemple, volem simplificar la següent fracció\begin{equation*}
\frac{x^{3}+5x^{2}+6x}{x^{3}+3x^{2}+2x}
\end{equation*}Descomponent en factors els polinomis numerador i denominador, obtenim\begin{equation*}
\frac{x^{3}+5x^{2}+6x}{x^{3}+3x^{2}+2x}=\frac{x(x+3)(x+2)}{x(x+1)(x+2)}=\frac{x+3}{x+1}
\end{equation*}Per tant,\begin{equation*}
\frac{x+3}{x+1}
\end{equation*}és la fracció irreduïble equivalent a la fracció dada.

\subsection{Operacions en el conjunt de les fraccions racionals}

Les definicions i regles per a efectuar la suma i el producte de dos
fraccions racionals són totalment anàlogues a les definicions i
regles per a les operacions amb nombres racionals.

\subsubsection{Suma de fraccions racionals}

Donades dues fraccions racionals, definim la fracció racional \textbf{suma} com la que té com a representant la suma dels representants de
les fraccions racionals donades. Per a efectuar la suma s'aplica la regla\begin{equation*}
\frac{p(x)}{q(x)}+\frac{r(x)}{s(x)}=\frac{p(x)\cdot s(x)+q(x)\cdot r(x)}{q(x)\cdot s(x)}
\end{equation*}Observa que aquesta fracció és algebraica, perquè, ni $q(x)$ ni $s(x)$ són
el polinomi zero, i, per tant, $q(x)\cdot s(x)$ tampoc ho és. A més,
és fàcil demostrar que aquesta suma no depèn dels representants
escollits i, en conseqüència, la operació está bien definida.

Les propietats associativa i commutativa, així com la existència del
element neutre i de la fracció racional oposada d’una qualsevol es
dedueix de la mateixa manera ja coneguda per als nombres racionals.
Recordem, simplement, que el element neutre és\begin{equation*}
\frac{0}{q(x)}
\end{equation*}amb el polinomi zero en el numerador i qualsevol polinomi no nul en el
denominador, i que la fracció racional oposada de\begin{equation*}
\frac{p(x)}{q(x)}
\end{equation*}és\begin{equation*}
\frac{-p(x)}{q(x)}
\end{equation*}

Per exemple, volem efectuar la següent suma de fraccions\begin{equation*}
\frac{x+3}{x^{2}+x}+\frac{x-2}{x^{2}+3x+2}
\end{equation*}Evidentment podem utilitzar la regla que hem donat al definir la suma de
fraccions. Tanmateix, resulta molt més còmode utilitzar el m.c.m.
dels polinomis denominadors i convertir les dues fraccions donades en
altres equivalents amb el mateix denominador; observa, perquè, que es
es procedeix de la mateixa manera que la suma de fraccions numèriques. Les
descomposicions factorials dels polinomis denominadors són\begin{equation*}
\begin{array}{l}
x^{2}+x=x(x+1) \\
x^{2}+3x+2=(x+1)(x+2)\end{array}\end{equation*}i, per tant, obtenim\begin{equation*}
\text{m.c.m.}\left( x^{2}+x,x^{2}+3x+2\right) =x(x+1)(x+2)
\end{equation*}Aleshores, podem escriure\begin{eqnarray*}
\frac{x+3}{x^{2}+x}+\frac{x-2}{x^{2}+3x+2} &=&\frac{(x+3)(x+2)}{x(x+1)(x+2)}+\frac{x(x-2)}{x(x+1)(x+2)} \\
&=&\frac{(x+3)(x+2)+x(x-2)}{x(x+1)(x+2)} \\
&=&\frac{2x^{2}+3x+6}{x^{3}+3x^{2}+2x}
\end{eqnarray*}

\subsubsection{Producte de fraccions racionals}

Donades dues fraccions racionals, definim la fracció racional \textbf{producte} com la que té com a representant el producte dels
representants de les fraccions racionals donades. Per a efectuar el producte
s'aplica la regla\begin{equation*}
\frac{p(x)}{q(x)}\cdot \frac{r(x)}{s(x)}=\frac{p(x)\cdot r(x)}{q(x)\cdot s(x)}
\end{equation*}Observa que aquesta fracció és algebraica, perquè, ni $q(x)$ ni $s(x)$ són
el polinomi zero, i, per tant, $q(x)\cdot s(x)$ tampoc ho és. A més,
és fàcil demostrar que aquest producte no depèn dels representants
escollits i, en conseqüència, la operació está bien definida.

Les propietats associativa i commutativa es demostren de la mateixa manera que
per als nombres racionals. L'element neutre està constituït per
la fracció racional que té com a representant qualsevol fracció
algebraica de numerador igual al denominador. La fracció racional
inversa de
\begin{equation*}
\frac{p(x)}{q(x)}
\end{equation*}és\begin{equation*}
\frac{q(x)}{p(x)}
\end{equation*}Observa que aquesta fracció serà algebraica si $p(x)$ no és el
polinomi nul. Per tant, totes les fraccions racionals tenen inversa excepte
la fracció racional zero\begin{equation*}
\frac{0}{q(x)}
\end{equation*}on $q(x)$ és qualsevol polinomi no nul. També es compleix la
propietat distributiva respecte de la suma.

Per exemple, volem trobar la fracció irreduïble equivalent al
següent producte de fraccions\begin{equation*}
\frac{1+x}{1-x}\cdot \frac{x}{1-x^{2}}
\end{equation*}Dit amb altres paraules, es tracta d'efectuar el producte i simplificar després la fracció resultant. Així, tenim\begin{eqnarray*}
\frac{1+x}{1-x}\cdot \frac{x}{1-x^{2}} &=&\frac{x(1+x)}{(1-x)(1-x^{2})} \\
&=&\frac{x(1+x)}{(1-x)^{2}(1+x)} \\
&=&\frac{x}{(1-x)^{2}} \\
&=&\frac{1}{1-2x+x^{2}}
\end{eqnarray*}

\subsection{El cos de les fraccions racionals}

Les propietats de la suma i el producte de fraccions racionals es
s'enuncien conjuntament dient que el conjunt de les fraccions racionals
amb la suma i el producte constitueixen un cos commutatiu, anomenat
\textbf{cos de les fraccions racionals} de $\mathbb{R}[x]$.

Ara ja podem donar resposta al problema plantejat a l'inici d'aquesta secció. Recordem que es tractava de resoldre l'equació següent\begin{equation*}
(x-1)\cdot A(x)=x^{2}+1
\end{equation*}

Evidentment podem escriure aquesta equació mitjançant fraccions
racionals de la manera següent\begin{equation*}
\frac{x-1}{1}\cdot \frac{A(x)}{1}=\frac{x^{2}+1}{1}
\end{equation*}on hem identificat cada polinomi $p(x)$ amb la fracció\begin{equation*}
\frac{p(x)}{1}
\end{equation*}Multiplicant ara els dos membres de l'equació per la fracció
racional inversa de
\begin{equation*}
\frac{x-1}{1}
\end{equation*}és a dir, multiplicant per\begin{equation*}
\frac{1}{x-1}
\end{equation*}obtenim\begin{eqnarray*}
\frac{1}{x-1}\cdot \left( \frac{x-1}{1}\cdot \frac{p(x)}{1}\right) &=&\frac{1}{x-1}\cdot \frac{x^{2}+1}{1} \\
\left( \frac{1}{x-1}\cdot \frac{x-1}{1}\right) \cdot \frac{p(x)}{1} &=&\frac{x^{2}+1}{x-1} \\
\frac{x-1}{x-1}\cdot \frac{p(x)}{1} &=&\frac{x^{2}+1}{x-1} \\
p(x) &=&\frac{x^{2}+1}{x-1}
\end{eqnarray*}

\section{Funcions polinòmiques}

S'anomena \textbf{funció polinòmica} de grau $n$ a tota funció
real de variable real $f$ definida per\begin{equation*}
f(x)=a_{0}+a_{1}x+a_{2}x^{2}+\cdots +a_{n-1}x^{n-1}+a_{n}x^{n}
\end{equation*}per a tot $x\in \mathbb{R}$, amb $a_{i}\in \mathbb{R}$ ($1\leq i\leq n$) i $a_{n}\neq 0$.

Observa que les operacions que hi cal fer per obtenir la imatge d'un
element $x$ es poden fer per a qualsevol nombre real i, per tant, una
funció polinòmica té com domini tot el conjunt dels nombres reals.

El valor que pren la funció per a $x=0$ coincideix amb el terme
independent del polinomi\begin{equation*}
f(0)=a_{0}
\end{equation*}i, per tant, la gràfica de la funció tallarà l'eix
d'ordenades en el punt $(0,a_{0})$. Els punts de tall de la gràfica amb
l'eix d'abscisses seran els punts $x$ tals que $f(x)=0$, és a dir,
els zeros de la funció; aquests punts es corresponen amb les arrels del polinomi.

A continuació tractarem alguns casos particulars de funcions polinòmiques així com les seves gràfiques.

\subsection{Funcions constants}

Una \textbf{funció constant} és una funció polinòmica de
grau zero. Tota funció constant $f$ ve definida per\begin{equation*}
f(x)=a
\end{equation*}per a tot $x\in \mathbb{R}$, amb $a\in \mathbb{R}$.

És fàcil veure que la gràfica d'aquesta funció és una recta
paral·lela a l'eix d'abscisses que passa pel punt $(0,a)$ i la qual equació és\begin{equation*}
y=a
\end{equation*}Per tant, el recorregut d'aquesta funció és redueix a un un sol punt $y=a$.

En la figura següent, hem dibuixat diferents funcions d'aquest tipus per a
diferents valors de $a$. La funció constant igual a $0$, o també anomenada \textbf{funció nula}, ve donada per $f(x)=0$ per a tot $x\in
\mathbb{R}$; la gràfica d'aquesta funció és la recta $y=0$ que
coincideix amb l'eix d’abscisses. \begin{center}
\includegraphics[width=2.5486in,height=2.5797in,keepaspectratio]{poli2.png}
\end{center}

\subsection{Funcions afins}

Una \textbf{funció afí} és una funció polinòmica de primer
grau. Tota funció afí $f$ ve definida per\begin{equation*}
f(x)=ax+b
\end{equation*}per a tot $x\in \mathbb{R}$, amb $a,b\in \mathbb{R}$ i $a\neq 0$.

És fàcil veure que la gràfica d'aquesta funció és una recta que
passa pels punts $(0,b)$ i $(-b/a,0)$, i la qual equació és\begin{equation*}
y=ax+b
\end{equation*}Per tant, el recorregut d'aquesta funció és tot el conjunt dels nombres reals.

Observa que si $a>0$, aleshores es compleix\begin{equation*}
\begin{array}{ccccc}
x_{1}\leq x_{2} & \Longrightarrow & ax_{1}\leq ax_{2} & \Longrightarrow &
f(x_{1})=ax_{1}+b\leq ax_{2}+b=f(x_{2})\end{array}\end{equation*}i, per tant, la funció és creixent. De la mateixa manera, si $a<0$,
aleshores es compleix\begin{equation*}
\begin{array}{ccccc}
x_{1}\leq x_{2} & \Longrightarrow & ax_{1}\geq ax_{2} & \Longrightarrow &
f(x_{1})=ax_{1}+b\geq ax_{2}+b=f(x_{2})\end{array}\end{equation*}i, per tant, la funció és decreixent.

Si $b=0$, la funció s’anomena \textbf{lineal} i la seva gràfica és una
recta de pendent $a$ que passa per l’origen de coordenades. En la
següent figura hem dibuixat les gràfiques d'algunes funcions d'
aquest tipus. \begin{center}
\includegraphics[width=2.674in,height=2.5581in,keepaspectratio]{poli3.png}
\end{center}

\subsection{Funcions quadràtiques}

Una \textbf{funció quadràtica} és una funció polinòmica de
segon grau. Tota funció quadràtica $f$ ve definida per\begin{equation*}
f(x)=ax^{2}+bx+c
\end{equation*}per a tot $x\in \mathbb{R}$, amb $a,b,c\in \mathbb{R}$ i $a\neq 0$.

Abans de passar a estudiar aquest cas general, tractarem alguns casos
especials més senzills.

\begin{enumerate}
\item $f(x)=x^{2}$
\end{enumerate}

És el cas més senzill de funció quadràtica. La gràfica d'
aquesta funció és representada per la paràbola $y=x^{2}$; hem
dibuixat aquesta corba en la figura següent \begin{center}
\includegraphics[width=2.6524in,height=2.6948in,keepaspectratio]{poli4.png}
\end{center}A partir de la gràfica d'aquesta
funció podem observar les següents propietats:

\begin{itemize}
\item La funció és sempre positiva

\item És simètrica respecte de l'eix d’ordenades, ja que $f(-x)=f(x)$; a l'eix de simetria s'anomena eix de la paràbola.

\item Presenta un punt d'ordenada mínima (mínim) en el punt $(0,0)$; a aquest punt s'anomena vèrtex de la paràbola.

\item Per a valors positius de la $x$ la funció és creixent, perquè\begin{equation*}
\begin{array}{ccc}
0<x_{1}\leq x_{2} & \Longrightarrow & f(x_{1})=x_{1}^{2}\leq
x_{2}^{2}=f(x_{2})\end{array}\end{equation*}i per a valors negatius de la $x$ la funció és decreixent, perquè\begin{equation*}
\begin{array}{ccc}
x_{1}\leq x_{2}<0 & \Longrightarrow & f(x_{1})=x_{1}^{2}\geq
x_{2}^{2}=f(x_{2})\end{array}\end{equation*}
\end{itemize}

\begin{enumerate}
\item[2.] $f(x)=-x^{2}$
\end{enumerate}

La gràfica d'aquesta funció s’obté per simetria respecte de l'
eix d'abscisses de la paràbola $y=x^{2}$. \begin{center}
\includegraphics[width=2.6636in,height=2.5071in,keepaspectratio]{poli5.png}
\end{center}En aquest cas la funció en lloc
de ser positiva sempre és negativa. La simetria respecte de l'eix
d'ordenades es conserva però canvia el creixement de la funció. Ara la
funció és creixent per a valors negatius de la $x$ i decreixent per a
valors positius. En lloc de tenir un punt d'ordenada mínima en $(0,0)$ ara té un punt d'ordenada màxima (màxim).

\begin{enumerate}
\item[3.] $f(x)=ax^{2}$, amb $a\in \mathbb{R}$ i $a\neq 0$
\end{enumerate}

La gràfica d'aquesta funció s'obtindrà a partir de la paràbola $y=x^{2}$ multiplicant les seves ordenades per la constant $a$. Serà
una paràbola amb les mateixes característiques que $y=x^{2}$ o $y=-x^{2}$, segons que $a$ sigui positiu o negatiu, respectivament. En la
figura següent hem dibuixat diverses paràboles del tipus $y=ax^{2}$; és
important observar que les branques de la paràbola seran més
tancades o obertes segons que $a$ augmenti o disminueixi en valor absolut.
\begin{center}
\includegraphics[width=2.9032in,height=2.8608in,keepaspectratio]{poli6.png}
\end{center}

\begin{enumerate}
\item[4.] $f(x)=ax^{2}+c$, amb $a,c\in \mathbb{R}$ i $a,c$ no nuls
\end{enumerate}

La gràfica d'aquesta funció s'obtindrà a partir de la paràbola $y=ax^{2}$ traslladant-la $c$ unitats seguint l'eix d'ordenades; es traslladarà $c$ unitats en el sentit positiu de l'eix d'ordenades si $c>0$, i $c$ unitats en el sentit negatiu d'aquest eix si $c<0$. Observa la figura següent per a adonar-te d'aquest fet. \begin{center}
\includegraphics[width=2.642in,height=2.5797in,keepaspectratio]{poli7.png}
\end{center}És
important observar que ara el vèrtex de la paràbola s'ha
traslladat del punt $(0,0)$ al punt $(0,c)$. També és important saber
que ara la funció no ha de ser necessàriament sempre positiva o sempre
negativa. En efecte, si $a>0$ i $c<0$, aleshores la paràbola $y=ax^{2}+c$
tallarà l'eix d'abscisses en els punts\begin{equation*}
\left( \sqrt{-\frac{c}{a}},0\right) \text{ \ \ \ i \ \ \ }\left( -\sqrt{-\frac{c}{a}},0\right)
\end{equation*}ja que\begin{equation*}
a\left( \pm \sqrt{-\frac{c}{a}}\right) ^{2}+c=a\left( -\frac{c}{a}\right)
+c=0
\end{equation*}Observa que passal mateix si $a<0$ i $c>0$. En resum, si $a$ i $c$
té signes oposats, la paràbola talla l'eix d'abscisses en dos
punts diferents i, com a conseqüència, la funció canvia de signe.

\begin{enumerate}
\item[5.] $f(x)=a(x-p)^{2}+q$, amb $a,p,q\in \mathbb{R}$ i $a,p,q$ no nuls
\end{enumerate}

La gràfica d'aquesta funció s'obtindrà a partir de la paràbola $y=ax^{2}+q$ traslladant-la $p$ unitats seguint l'eix d'abscisses; es traslladarà $p$ unitats en el sentit positiu de l'eix d'abscisses si $p>0$, i $p$ unitats en el sentit negatiu d'aquest eix si $p<0 $. Observa la figura següent per a adonar-te d'aquest fet. \begin{center}
\includegraphics[width=2.6636in,height=2.6636in,keepaspectratio]{poli8.png}
\end{center}

És important observar que ara el vèrtex de la paràbola s'ha
traslladat del punt $(0,q)$ al punt $(p,q)$. L'eix de simetria de la
paràbola és ara paral·lel a l'eix d’ordenades i passa pel punt $(p,q)$. És fàcil comprovar que per a $a>0$, la funció és
decreixent a l'esquerra del vèrtex i creixent a la seva dreta, i la seva gràfica presenta un mínim en el vèrtex; anàlogament, es
comprova que si $a<0$, la funció és creixent a l'esquerra del vèrtex i decreixent a la seva dreta, i la seva gràfica presenta un màxim en
el vèrtex.

\begin{enumerate}
\item[6.] $f(x)=ax^{2}+bx+c$, amb $a,b,c\in \mathbb{R}$ i $a\neq 0$
\end{enumerate}

Considerem finalment el cas general. Observa en primer lloc que podem
escriure\begin{equation*}
ax^{2}+bx+c=a\left( x+\frac{b}{2a}\right) ^{2}+\left( c-\frac{b^{2}}{4a}\right)
\end{equation*}Per exemple,\begin{equation*}
2x^{2}+16x+34=2(x+4)^{2}+2
\end{equation*}Per tant, si fem\begin{eqnarray*}
p &=&\frac{b}{2a} \\
q &=&c-\frac{b^{2}}{4a}
\end{eqnarray*}aleshores\begin{equation*}
ax^{2}+bx+c=a(x-p)^{2}+q
\end{equation*}de manera que la gràfica d'aquesta funció és una paràbola de la
forma\begin{equation*}
y=a(x-p)^{2}+q
\end{equation*}Com ja hem vist les característiques d'aquesta paràbola són:

\begin{itemize}
\item Vèrtex en el punt $\left( -\frac{b}{2a},c-\frac{b^{2}}{4a}\right) $

\item Eix de simetria paral·lel a l'eix d’ordenades i passa pel vèrtex

\item Per a $a>0$, la funció presenta un punt d'ordenada mínima
en el vèrtex (mínim) i, per a $a<0$, la funció presenta un
punt d'ordenada màxima en el vèrtex (màxim). Per tant, si $a>0$, el recorregut de la funció és $\left[ c-\frac{b^{2}}{4a},+\infty
\right) $, i si $a<0$, és $\left( -\infty,c-\frac{b^{2}}{4a}\right] $

\item Per a $a>0$, diem que la paràbola té concavitat positiva
(convexa) i, per a $a<0$, diem que la paràbola té concavitat
negativa (còncava)

\item La paràbola talla l'eix d’ordenades en el punt $(0,c)$; talla
al eix d'abscisses si $a$ i $c-\frac{b^{2}}{4a}$ té signes oposats o $c-\frac{b^{2}}{4a}=0$ (la paràbola és aleshores tangent a l'eix d’abscisses
y el punt de contacte és el vèrtex). Observa que perquè la paràbola talli l'eix d'abscisses el polinomi $ax^{2}+bx+c$ ha de tenir arrels reals i, per tant, si $a>0$ la condició perquè això passi és\begin{equation*}
c-\frac{b^{2}}{4a}\leq 0
\end{equation*}equivalent a\begin{equation*}
b^{2}-4ac\geq 0
\end{equation*}Com a conseqüència, el polinomi $ax^{2}+bx+c$ no tindrà arrels
reals si $b^{2}-4ac<0$.
\end{itemize}

Per exemple, en la figura següent hem representat la gràfica de la
funció $f(x)=x^{2}-2x-3$. \begin{center}
\includegraphics[width=2.5694in,height=2.5789in,keepaspectratio]{poli9.png}
\end{center}Com que\begin{equation*}
x^{2}-2x-3=(x-1)^{2}-4
\end{equation*}aleshores $(1,-4)$ és el vèrtex de la paràbola. El seu eix de simetria és la recta vertical d'equació $x=1$. La funció té un mínim en $x=1$ el valor del qual és $-4$. La gràfica és convexa i talla l’
eix d’ordenades en el punt $(0,-3)$. Els punts de tall amb l'eix d'abscisses són les arrels del polinomi $x^{2}-2x-3$. Com que\begin{equation*}
x^{2}-2x-3=(x+1)(x-3)
\end{equation*}aleshores les arrels són $-1$ i $3$, i, per tant, els punts de tall
amb l'eix d'abscisses són $(-1,0)$ i $(3,0)$.

\subsection{Altres funcions polinòmiques senzilles}

En l'apartat anterior s'ha vist com a partir de la paràbola $y=x^{2}$ es pot representar qualsevol altra paràbola de la forma $y=ax^{2}+bx+c$, mitjançant simetries o translacions dels eixos. La mateixa
tècnica pot servir-nos per a representar altres funcions polinòmiques
senzilles.

Per exemple, a partir de la gràfica de la funció $f(x)=x^{3}$,
obtinguda mitjançant una taula de valors, \begin{center}
\includegraphics[width=2.6109in,height=2.5382in,keepaspectratio]{poli10.png}
\end{center}podem obtenir altres gràfiques de
funcions polinòmiques de tercer grau.

Per exemple, la gràfica de la funció\begin{equation*}
f(x)=x^{3}-3x^{2}+3x+3
\end{equation*}pot obtenir-se a partir de $y=x^{3}$, perquè\begin{equation*}
x^{3}-3x^{2}+3x+3=(x-1)^{3}+4
\end{equation*}En efecte, de la gràfica de $y=x^{3}$, mitjançant una translació
vertical, seguint l'eix d’ordenades en el sentit positiu $4$ unitats,
obtindrem la gràfica de $y=x^{3}+4$. D'aquesta gràfica, mitjançant una
translació horizontal, seguint l'eix d'abscisses en el sentit
positiu $1$ unitat, obtindrem la gràfica de $y=(x-1)^{3}+4$, \begin{center}
\includegraphics[width=2.5797in,height=2.5382in,keepaspectratio]{poli11.png}
\end{center}y, a partir de la seva gràfica, podem deduir algunes de les seves
propietats.

Evidentment, a partir de la gràfica de la funció $f(x)=x^{4}$,
obtinguda mitjançant una taula de valors, \begin{center}
\includegraphics[width=2.725in,height=2.6844in,keepaspectratio]{poli12.png}
\end{center}podem obtenir altres gràfiques de
funcions polinòmiques de quart grau.

Per exemple, la gràfica de la funció\begin{equation*}
f(x)=-x^{4}-4x^{3}-6x^{3}-4x+4
\end{equation*}pot obtenir-se a partir de $y=x^{4}$, perquè\begin{equation*}
-x^{4}-4x^{3}-6x^{2}-4x+4=-(x+1)^{4}+5
\end{equation*}Recorda que\begin{equation*}
(x+1)^{4}=x^{4}+4x^{3}+6x^{2}+4x+1
\end{equation*}En efecte, de la gràfica de $y=x^{4}$, mitjançant una simetria
respecte de l'eix d’abscisses, obtindrem la gràfica de $y=-x^{4}$.
Ara, mitjançant una translació vertical, seguint l'eix d’ordenades
en el sentit positiu $5$ unitats, obtindrem la gràfica de $y=-x^{4}+5$. D'aquesta gràfica, mitjançant una translació horizontal,
seguint l'eix d'abscisses en el sentit negatiu $1$ unitat, obtindrem
la gràfica de $y=-(x+1)^{4}+5$, \begin{center}
\includegraphics[width=2.7882in,height=2.5797in,keepaspectratio]{poli13.png}
\end{center}y, a partir de la seva gràfica,
podem deduir algunes de les seves propietats.

\section{Funcions racionals}

Així com cada polinomi dona lloc a una funció polinòmica, cada
fracció algebraica dona lloc a una funció que s’anomena \textbf{funció racional}, associada a una fracció algebraica. Una funció
racional és tota funció real de variable real $f$ definida per\begin{equation*}
f(x)=\frac{p(x)}{q(x)}
\end{equation*}per a tot $x\in \mathbb{R}$, on $p(x),q(x)\in \mathbb{R}[x]$ i $q(x)$ no
és el polinomi zero.

Per exemple,\begin{equation*}
f(x)=\frac{x^{2}-1}{3x-2}
\end{equation*}és una funció racional. Encara que no tractarem aquí aquest tipus de
funcions, només farem algunes observacions. En buscar la imatge d'un
nombre real per $f$ pot passar que per a aquest valor s'anul·li el
denominador, cas en què el quocient no té sentit. Aquests valors de $x$
que anul·len el denominador no tenen imatge i cal prescindir-ne. Per tant, el domini d’una funció racional és\begin{eqnarray*}
\text{Dom }f &=&\left\{ x\in \mathbb{R}:q(x)\neq 0\right\} \\
&=&\mathbb{R}-\left\{ x\in \mathbb{R}:q(x)=0\right\}
\end{eqnarray*}

En l'exemple anterior, el valor $x=\frac{2}{3}$ no té imatge i, per
tant,\begin{eqnarray*}
\text{Dom }f &=&\mathbb{R}-\left\{ x\in \mathbb{R}:3x-2=0\right\} \\
&=&\mathbb{R}-\left\{ \frac{2}{3}\right\}
\end{eqnarray*}

Si s'efectua la representació gràfica d’una funció
racional caldrà prendre la precaució de assenyalar els punts $x$
que no tenen imatge. De moment, per a efectuar la representació gràfica de la funció\begin{equation*}
f(x)=\frac{x^{2}-1}{3x-2}
\end{equation*}podem utilitzar una taula de valors i assenyalar el valor $x=\frac{2}{3}$. Cap valor de $y$ associat a $\frac{2}{3}$ pot ser punt de la gràfica, és a dir, no és possible que la recta vertical $x=\frac{2}{3}$
talli la gràfica de la funció. Com que el valor $\frac{2}{3}$
no es pot donar, tractarem de completar la taula amb valors pròxims a
aquest punt. En la figura següent es mostra la gràfica d'aquesta funció\begin{center}
\includegraphics[width=2.6316in,height=2.6109in,keepaspectratio]{poli14.png}
\end{center}En aquesta gràfica hem dibuixat també la recta
\begin{equation*}
y=\frac{1}{3}x+\frac{2}{9}
\end{equation*}a la que les branques de la gràfica s'aproximen a mesura que $x$ augmenta o
disminueix indefinidament. Veurem en un altre apartat el significat d'aquesta
recta així com el de la recta\begin{equation*}
x=\frac{2}{3}
\end{equation*}
